% =====================================================================
%  Fragmento integrable en main.tex mediante:  % =====================================================================
%  Fragmento integrable en main.tex mediante:  % =====================================================================
%  Fragmento integrable en main.tex mediante:  % =====================================================================
%  Fragmento integrable en main.tex mediante:  \input{generated/isoeficiencia.tex}
%  Reutiliza la notación, los paquetes (amsmath, booktabs, float) y las
%  complejidades ya derivadas en el documento principal. No define preámbulo.
% =====================================================================

\subsection{Análisis de isoeficiencia (condición de escalabilidad)}

En este curso la \textbf{isoeficiencia} se evalúa como una \emph{condición de escalabilidad
débil} (weak scaling) basada en la Ley de Gustafson. No se busca una función de isoeficiencia
aislada, sino la relación que dicta cómo debe crecer el tamaño del problema $n$ en función del
número de procesos $p$ para que la eficiencia $E$ se mantenga constante ($E=O(1)$). Formalizamos
aquí la condición que ya se anticipó de manera cualitativa en las secciones de escalabilidad
fuerte y débil, y comparamos el paradigma usado (MPI) con el de memoria compartida (OMP/PRAM).

\subsubsection{Función de overhead y forma canónica de la eficiencia}

Para obtener una única condición de escalabilidad tomamos un tamaño de problema $n$ con
$n_{tr}=n_{te}=n$ (el dataset completo escala). De las complejidades ya derivadas, el trabajo
total (tiempo secuencial dominante, omitiendo $\log K$) es
\[
W \equiv T_1 = O(n^2 d),
\]
y el cómputo paralelo reparte ese trabajo idealmente, $T_{\text{comp}} = O(n^2 d / p)$.

Definimos la \textbf{función de overhead total} $T_o$ como el trabajo que no realiza el algoritmo
secuencial:
\[
T_o(p) = p\,T_p - W = p\,T_{\text{comp}} + p\,T_{\text{comm}} - W.
\]
Como $p\,T_{\text{comp}} = W$, todo el overhead proviene de la comunicación y la sincronización:
\[
\boxed{\; T_o(p) = p\, T_{\text{comm}}. \;}
\]
Sustituyendo en la eficiencia se obtiene la forma canónica del curso:
\[
E = \frac{S_p}{p} = \frac{T_1}{p\,T_p} = \frac{W}{W + T_o}
  = \frac{1}{\,1 + \dfrac{T_o}{W}\,}.
\]
Para que $E=O(1)$ sea constante, el término de overhead $T_o/W$ debe permanecer constante: esa
igualdad, despejada, define la condición de isoeficiencia.

\subsubsection{Isoeficiencia en MPI (memoria distribuida)}

En MPI la comunicación es explícita y es el cuello de botella. Sustituyendo $n_{tr}=n_{te}=n$ en
$T_{\text{comm}}$ y conservando los términos dominantes:
\[
T_o = p\,T_{\text{comm}}
    \approx \underbrace{p^2\alpha}_{(1)\ \text{latencia}}
    \;+\; \underbrace{p\,\beta\, n\, d}_{(2)\ \text{scatter}}
    \;+\; \underbrace{p\,\beta\, n\, d\,\log p}_{(3)\ \text{bcast/reducción}}.
\]
El término (3) proviene de difundir el conjunto de prueba ($O(nd)$ bytes) y de reducir las listas
Top-$K$ a lo largo de una jerarquía de profundidad $\log p$; domina sobre (2). Imponemos
$T_o/W=\text{const}$ con $W=n^2 d$:
\[
(1)\quad \frac{p^2\alpha}{n^2 d}=\text{const} \;\Rightarrow\; n \propto p,
\qquad
(2)\quad \frac{p\,\beta n d}{n^2 d}=\frac{p\beta}{n}=\text{const} \;\Rightarrow\; n \propto p,
\]
\[
(3)\quad \frac{p\,\beta n d\,\log p}{n^2 d}=\frac{p\beta\log p}{n}=\text{const}
\;\Rightarrow\; n \propto p\,\log p.
\]
La condición más restrictiva es la del término (3). Por tanto:
\[
\boxed{\;
\text{Isoeficiencia MPI:}\quad n \propto p\,\log p,
\qquad
W_{iso} = \Theta\!\big(p^2 \log^2 p\big).
\;}
\]
Es decir, el dataset debe crecer ligeramente más rápido que $p$ (factor $\log p$) para amortizar
la difusión del test y la reducción global de candidatos sobre la red. Si la reducción se hiciera
con un \texttt{gather} plano al proceso raíz en vez de la reducción en árbol, el término dominante
sería $O(p)$ en lugar de $O(\log p)$, saturando el enlace del root y empeorando notablemente la
isoeficiencia.

\subsubsection{Isoeficiencia en OMP / PRAM (memoria compartida)}

En OpenMP los hilos comparten el mismo espacio de direcciones (modelo UMA), lo que cambia el
overhead de forma fundamental:
\begin{itemize}
    \item \textbf{No hay difusión de datos:} el conjunto de prueba es visible por todos los hilos,
    por lo que \emph{desaparecen} los términos $p\alpha$ y $\beta n d$ de la comunicación
    explícita ($T_{\text{comm}} \approx 0$).
    \item El único overhead proviene de la \textbf{sincronización de hilos} (barreras) y de la
    reducción de listas Top-$K$ de profundidad $O(\log p)$, realizada mediante accesos directos a
    memoria.
\end{itemize}
Modelando ese overhead (con $K$ constante):
\[
T_o^{\text{omp}} \approx p\cdot O(n_{te}K\log p) + O(p)
= O(p\,n\,\log p) + O(p).
\]
Imponiendo $T_o^{\text{omp}}/W=\text{const}$ con $W=n^2 d$, y observando que ya no aparece el
factor de transporte de datos $\beta d$:
\[
\frac{p\,n\,\log p}{n^2 d}=\frac{p\log p}{n\,d}=\text{const}
\;\Rightarrow\;
n \propto \frac{p\,\log p}{d}.
\]
\[
\boxed{\;
\text{Isoeficiencia OMP:}\quad n \propto \frac{p\,\log p}{d}\ \ (\text{a lo sumo } n\propto p),
\qquad
W_{iso} = \Theta(p^2)\ \text{sin la constante de red } \beta.
\;}
\]
OMP tiene una isoeficiencia asintóticamente mejor (no paga latencia ni ancho de banda de red),
\textbf{pero} $p$ está estrictamente acotado al número de núcleos físicos del nodo y todo el
dataset debe caber en su RAM.

\subsubsection{Comparación y elección del paradigma}

\begin{table}[H]
\centering
\begin{tabular}{lll}
\toprule
Criterio & MPI (memoria distribuida) & OMP (memoria compartida) \\
\midrule
Comunicación $T_{\text{comm}}$ & Explícita: $\alpha,\beta$ (red) & $\approx 0$ (acceso directo) \\
Overhead dominante & Difusión/reducción $\beta n d\log p$ & Sincronización (barreras) \\
Condición de isoeficiencia & $n \propto p\log p$ & $n \propto p\log p / d$ (más suave) \\
Trabajo $W_{iso}$ & $\Theta(p^2\log^2 p)$ & $\Theta(p^2)$ sin $\beta$ \\
Límite de $p$ & Horizontal (varios nodos) & Núcleos de un solo nodo \\
Límite de $n$ & Suma de RAM de los nodos & RAM de un solo nodo \\
\bottomrule
\end{tabular}
\caption{Comparación de la escalabilidad de ambos paradigmas para el KNN.}
\label{tab:iso-comparacion}
\end{table}

Aunque OMP escala mejor asintóticamente, su paralelismo está limitado a los núcleos y la memoria
de un solo nodo. En este proyecto se eligió \textbf{MPI por un tema de memoria}: al distribuir el
conjunto de entrenamiento, el mismo código puede superar el límite de RAM de un nodo y escalar
$n$ y $p$ horizontalmente sin reescritura. Los experimentos evaluaron hasta $p=32$ procesos, y la
elección del paradigma de memoria distribuida es la que habilita crecer más allá de ese punto
cuando el dataset no cabe en un único nodo.


%  Reutiliza la notación, los paquetes (amsmath, booktabs, float) y las
%  complejidades ya derivadas en el documento principal. No define preámbulo.
% =====================================================================

\subsection{Análisis de isoeficiencia (condición de escalabilidad)}

En este curso la \textbf{isoeficiencia} se evalúa como una \emph{condición de escalabilidad
débil} (weak scaling) basada en la Ley de Gustafson. No se busca una función de isoeficiencia
aislada, sino la relación que dicta cómo debe crecer el tamaño del problema $n$ en función del
número de procesos $p$ para que la eficiencia $E$ se mantenga constante ($E=O(1)$). Formalizamos
aquí la condición que ya se anticipó de manera cualitativa en las secciones de escalabilidad
fuerte y débil, y comparamos el paradigma usado (MPI) con el de memoria compartida (OMP/PRAM).

\subsubsection{Función de overhead y forma canónica de la eficiencia}

Para obtener una única condición de escalabilidad tomamos un tamaño de problema $n$ con
$n_{tr}=n_{te}=n$ (el dataset completo escala). De las complejidades ya derivadas, el trabajo
total (tiempo secuencial dominante, omitiendo $\log K$) es
\[
W \equiv T_1 = O(n^2 d),
\]
y el cómputo paralelo reparte ese trabajo idealmente, $T_{\text{comp}} = O(n^2 d / p)$.

Definimos la \textbf{función de overhead total} $T_o$ como el trabajo que no realiza el algoritmo
secuencial:
\[
T_o(p) = p\,T_p - W = p\,T_{\text{comp}} + p\,T_{\text{comm}} - W.
\]
Como $p\,T_{\text{comp}} = W$, todo el overhead proviene de la comunicación y la sincronización:
\[
\boxed{\; T_o(p) = p\, T_{\text{comm}}. \;}
\]
Sustituyendo en la eficiencia se obtiene la forma canónica del curso:
\[
E = \frac{S_p}{p} = \frac{T_1}{p\,T_p} = \frac{W}{W + T_o}
  = \frac{1}{\,1 + \dfrac{T_o}{W}\,}.
\]
Para que $E=O(1)$ sea constante, el término de overhead $T_o/W$ debe permanecer constante: esa
igualdad, despejada, define la condición de isoeficiencia.

\subsubsection{Isoeficiencia en MPI (memoria distribuida)}

En MPI la comunicación es explícita y es el cuello de botella. Sustituyendo $n_{tr}=n_{te}=n$ en
$T_{\text{comm}}$ y conservando los términos dominantes:
\[
T_o = p\,T_{\text{comm}}
    \approx \underbrace{p^2\alpha}_{(1)\ \text{latencia}}
    \;+\; \underbrace{p\,\beta\, n\, d}_{(2)\ \text{scatter}}
    \;+\; \underbrace{p\,\beta\, n\, d\,\log p}_{(3)\ \text{bcast/reducción}}.
\]
El término (3) proviene de difundir el conjunto de prueba ($O(nd)$ bytes) y de reducir las listas
Top-$K$ a lo largo de una jerarquía de profundidad $\log p$; domina sobre (2). Imponemos
$T_o/W=\text{const}$ con $W=n^2 d$:
\[
(1)\quad \frac{p^2\alpha}{n^2 d}=\text{const} \;\Rightarrow\; n \propto p,
\qquad
(2)\quad \frac{p\,\beta n d}{n^2 d}=\frac{p\beta}{n}=\text{const} \;\Rightarrow\; n \propto p,
\]
\[
(3)\quad \frac{p\,\beta n d\,\log p}{n^2 d}=\frac{p\beta\log p}{n}=\text{const}
\;\Rightarrow\; n \propto p\,\log p.
\]
La condición más restrictiva es la del término (3). Por tanto:
\[
\boxed{\;
\text{Isoeficiencia MPI:}\quad n \propto p\,\log p,
\qquad
W_{iso} = \Theta\!\big(p^2 \log^2 p\big).
\;}
\]
Es decir, el dataset debe crecer ligeramente más rápido que $p$ (factor $\log p$) para amortizar
la difusión del test y la reducción global de candidatos sobre la red. Si la reducción se hiciera
con un \texttt{gather} plano al proceso raíz en vez de la reducción en árbol, el término dominante
sería $O(p)$ en lugar de $O(\log p)$, saturando el enlace del root y empeorando notablemente la
isoeficiencia.

\subsubsection{Isoeficiencia en OMP / PRAM (memoria compartida)}

En OpenMP los hilos comparten el mismo espacio de direcciones (modelo UMA), lo que cambia el
overhead de forma fundamental:
\begin{itemize}
    \item \textbf{No hay difusión de datos:} el conjunto de prueba es visible por todos los hilos,
    por lo que \emph{desaparecen} los términos $p\alpha$ y $\beta n d$ de la comunicación
    explícita ($T_{\text{comm}} \approx 0$).
    \item El único overhead proviene de la \textbf{sincronización de hilos} (barreras) y de la
    reducción de listas Top-$K$ de profundidad $O(\log p)$, realizada mediante accesos directos a
    memoria.
\end{itemize}
Modelando ese overhead (con $K$ constante):
\[
T_o^{\text{omp}} \approx p\cdot O(n_{te}K\log p) + O(p)
= O(p\,n\,\log p) + O(p).
\]
Imponiendo $T_o^{\text{omp}}/W=\text{const}$ con $W=n^2 d$, y observando que ya no aparece el
factor de transporte de datos $\beta d$:
\[
\frac{p\,n\,\log p}{n^2 d}=\frac{p\log p}{n\,d}=\text{const}
\;\Rightarrow\;
n \propto \frac{p\,\log p}{d}.
\]
\[
\boxed{\;
\text{Isoeficiencia OMP:}\quad n \propto \frac{p\,\log p}{d}\ \ (\text{a lo sumo } n\propto p),
\qquad
W_{iso} = \Theta(p^2)\ \text{sin la constante de red } \beta.
\;}
\]
OMP tiene una isoeficiencia asintóticamente mejor (no paga latencia ni ancho de banda de red),
\textbf{pero} $p$ está estrictamente acotado al número de núcleos físicos del nodo y todo el
dataset debe caber en su RAM.

\subsubsection{Comparación y elección del paradigma}

\begin{table}[H]
\centering
\begin{tabular}{lll}
\toprule
Criterio & MPI (memoria distribuida) & OMP (memoria compartida) \\
\midrule
Comunicación $T_{\text{comm}}$ & Explícita: $\alpha,\beta$ (red) & $\approx 0$ (acceso directo) \\
Overhead dominante & Difusión/reducción $\beta n d\log p$ & Sincronización (barreras) \\
Condición de isoeficiencia & $n \propto p\log p$ & $n \propto p\log p / d$ (más suave) \\
Trabajo $W_{iso}$ & $\Theta(p^2\log^2 p)$ & $\Theta(p^2)$ sin $\beta$ \\
Límite de $p$ & Horizontal (varios nodos) & Núcleos de un solo nodo \\
Límite de $n$ & Suma de RAM de los nodos & RAM de un solo nodo \\
\bottomrule
\end{tabular}
\caption{Comparación de la escalabilidad de ambos paradigmas para el KNN.}
\label{tab:iso-comparacion}
\end{table}

Aunque OMP escala mejor asintóticamente, su paralelismo está limitado a los núcleos y la memoria
de un solo nodo. En este proyecto se eligió \textbf{MPI por un tema de memoria}: al distribuir el
conjunto de entrenamiento, el mismo código puede superar el límite de RAM de un nodo y escalar
$n$ y $p$ horizontalmente sin reescritura. Los experimentos evaluaron hasta $p=32$ procesos, y la
elección del paradigma de memoria distribuida es la que habilita crecer más allá de ese punto
cuando el dataset no cabe en un único nodo.


%  Reutiliza la notación, los paquetes (amsmath, booktabs, float) y las
%  complejidades ya derivadas en el documento principal. No define preámbulo.
% =====================================================================

\subsection{Análisis de isoeficiencia (condición de escalabilidad)}

En este curso la \textbf{isoeficiencia} se evalúa como una \emph{condición de escalabilidad
débil} (weak scaling) basada en la Ley de Gustafson. No se busca una función de isoeficiencia
aislada, sino la relación que dicta cómo debe crecer el tamaño del problema $n$ en función del
número de procesos $p$ para que la eficiencia $E$ se mantenga constante ($E=O(1)$). Formalizamos
aquí la condición que ya se anticipó de manera cualitativa en las secciones de escalabilidad
fuerte y débil, y comparamos el paradigma usado (MPI) con el de memoria compartida (OMP/PRAM).

\subsubsection{Función de overhead y forma canónica de la eficiencia}

Para obtener una única condición de escalabilidad tomamos un tamaño de problema $n$ con
$n_{tr}=n_{te}=n$ (el dataset completo escala). De las complejidades ya derivadas, el trabajo
total (tiempo secuencial dominante, omitiendo $\log K$) es
\[
W \equiv T_1 = O(n^2 d),
\]
y el cómputo paralelo reparte ese trabajo idealmente, $T_{\text{comp}} = O(n^2 d / p)$.

Definimos la \textbf{función de overhead total} $T_o$ como el trabajo que no realiza el algoritmo
secuencial:
\[
T_o(p) = p\,T_p - W = p\,T_{\text{comp}} + p\,T_{\text{comm}} - W.
\]
Como $p\,T_{\text{comp}} = W$, todo el overhead proviene de la comunicación y la sincronización:
\[
\boxed{\; T_o(p) = p\, T_{\text{comm}}. \;}
\]
Sustituyendo en la eficiencia se obtiene la forma canónica del curso:
\[
E = \frac{S_p}{p} = \frac{T_1}{p\,T_p} = \frac{W}{W + T_o}
  = \frac{1}{\,1 + \dfrac{T_o}{W}\,}.
\]
Para que $E=O(1)$ sea constante, el término de overhead $T_o/W$ debe permanecer constante: esa
igualdad, despejada, define la condición de isoeficiencia.

\subsubsection{Isoeficiencia en MPI (memoria distribuida)}

En MPI la comunicación es explícita y es el cuello de botella. Sustituyendo $n_{tr}=n_{te}=n$ en
$T_{\text{comm}}$ y conservando los términos dominantes:
\[
T_o = p\,T_{\text{comm}}
    \approx \underbrace{p^2\alpha}_{(1)\ \text{latencia}}
    \;+\; \underbrace{p\,\beta\, n\, d}_{(2)\ \text{scatter}}
    \;+\; \underbrace{p\,\beta\, n\, d\,\log p}_{(3)\ \text{bcast/reducción}}.
\]
El término (3) proviene de difundir el conjunto de prueba ($O(nd)$ bytes) y de reducir las listas
Top-$K$ a lo largo de una jerarquía de profundidad $\log p$; domina sobre (2). Imponemos
$T_o/W=\text{const}$ con $W=n^2 d$:
\[
(1)\quad \frac{p^2\alpha}{n^2 d}=\text{const} \;\Rightarrow\; n \propto p,
\qquad
(2)\quad \frac{p\,\beta n d}{n^2 d}=\frac{p\beta}{n}=\text{const} \;\Rightarrow\; n \propto p,
\]
\[
(3)\quad \frac{p\,\beta n d\,\log p}{n^2 d}=\frac{p\beta\log p}{n}=\text{const}
\;\Rightarrow\; n \propto p\,\log p.
\]
La condición más restrictiva es la del término (3). Por tanto:
\[
\boxed{\;
\text{Isoeficiencia MPI:}\quad n \propto p\,\log p,
\qquad
W_{iso} = \Theta\!\big(p^2 \log^2 p\big).
\;}
\]
Es decir, el dataset debe crecer ligeramente más rápido que $p$ (factor $\log p$) para amortizar
la difusión del test y la reducción global de candidatos sobre la red. Si la reducción se hiciera
con un \texttt{gather} plano al proceso raíz en vez de la reducción en árbol, el término dominante
sería $O(p)$ en lugar de $O(\log p)$, saturando el enlace del root y empeorando notablemente la
isoeficiencia.

\subsubsection{Isoeficiencia en OMP / PRAM (memoria compartida)}

En OpenMP los hilos comparten el mismo espacio de direcciones (modelo UMA), lo que cambia el
overhead de forma fundamental:
\begin{itemize}
    \item \textbf{No hay difusión de datos:} el conjunto de prueba es visible por todos los hilos,
    por lo que \emph{desaparecen} los términos $p\alpha$ y $\beta n d$ de la comunicación
    explícita ($T_{\text{comm}} \approx 0$).
    \item El único overhead proviene de la \textbf{sincronización de hilos} (barreras) y de la
    reducción de listas Top-$K$ de profundidad $O(\log p)$, realizada mediante accesos directos a
    memoria.
\end{itemize}
Modelando ese overhead (con $K$ constante):
\[
T_o^{\text{omp}} \approx p\cdot O(n_{te}K\log p) + O(p)
= O(p\,n\,\log p) + O(p).
\]
Imponiendo $T_o^{\text{omp}}/W=\text{const}$ con $W=n^2 d$, y observando que ya no aparece el
factor de transporte de datos $\beta d$:
\[
\frac{p\,n\,\log p}{n^2 d}=\frac{p\log p}{n\,d}=\text{const}
\;\Rightarrow\;
n \propto \frac{p\,\log p}{d}.
\]
\[
\boxed{\;
\text{Isoeficiencia OMP:}\quad n \propto \frac{p\,\log p}{d}\ \ (\text{a lo sumo } n\propto p),
\qquad
W_{iso} = \Theta(p^2)\ \text{sin la constante de red } \beta.
\;}
\]
OMP tiene una isoeficiencia asintóticamente mejor (no paga latencia ni ancho de banda de red),
\textbf{pero} $p$ está estrictamente acotado al número de núcleos físicos del nodo y todo el
dataset debe caber en su RAM.

\subsubsection{Comparación y elección del paradigma}

\begin{table}[H]
\centering
\begin{tabular}{lll}
\toprule
Criterio & MPI (memoria distribuida) & OMP (memoria compartida) \\
\midrule
Comunicación $T_{\text{comm}}$ & Explícita: $\alpha,\beta$ (red) & $\approx 0$ (acceso directo) \\
Overhead dominante & Difusión/reducción $\beta n d\log p$ & Sincronización (barreras) \\
Condición de isoeficiencia & $n \propto p\log p$ & $n \propto p\log p / d$ (más suave) \\
Trabajo $W_{iso}$ & $\Theta(p^2\log^2 p)$ & $\Theta(p^2)$ sin $\beta$ \\
Límite de $p$ & Horizontal (varios nodos) & Núcleos de un solo nodo \\
Límite de $n$ & Suma de RAM de los nodos & RAM de un solo nodo \\
\bottomrule
\end{tabular}
\caption{Comparación de la escalabilidad de ambos paradigmas para el KNN.}
\label{tab:iso-comparacion}
\end{table}

Aunque OMP escala mejor asintóticamente, su paralelismo está limitado a los núcleos y la memoria
de un solo nodo. En este proyecto se eligió \textbf{MPI por un tema de memoria}: al distribuir el
conjunto de entrenamiento, el mismo código puede superar el límite de RAM de un nodo y escalar
$n$ y $p$ horizontalmente sin reescritura. Los experimentos evaluaron hasta $p=32$ procesos, y la
elección del paradigma de memoria distribuida es la que habilita crecer más allá de ese punto
cuando el dataset no cabe en un único nodo.


%  Reutiliza la notación, los paquetes (amsmath, booktabs, float) y las
%  complejidades ya derivadas en el documento principal. No define preámbulo.
% =====================================================================

\subsection{Análisis de isoeficiencia (condición de escalabilidad)}

En este curso la \textbf{isoeficiencia} se evalúa como una \emph{condición de escalabilidad
débil} (weak scaling) basada en la Ley de Gustafson. No se busca una función de isoeficiencia
aislada, sino la relación que dicta cómo debe crecer el tamaño del problema $n$ en función del
número de procesos $p$ para que la eficiencia $E$ se mantenga constante ($E=O(1)$). Formalizamos
aquí la condición que ya se anticipó de manera cualitativa en las secciones de escalabilidad
fuerte y débil, y comparamos el paradigma usado (MPI) con el de memoria compartida (OMP/PRAM).

\subsubsection{Función de overhead y forma canónica de la eficiencia}

Para obtener una única condición de escalabilidad tomamos un tamaño de problema $n$ con
$n_{tr}=n_{te}=n$ (el dataset completo escala). De las complejidades ya derivadas, el trabajo
total (tiempo secuencial dominante, omitiendo $\log K$) es
\[
W \equiv T_1 = O(n^2 d),
\]
y el cómputo paralelo reparte ese trabajo idealmente, $T_{\text{comp}} = O(n^2 d / p)$.

Definimos la \textbf{función de overhead total} $T_o$ como el trabajo que no realiza el algoritmo
secuencial:
\[
T_o(p) = p\,T_p - W = p\,T_{\text{comp}} + p\,T_{\text{comm}} - W.
\]
Como $p\,T_{\text{comp}} = W$, todo el overhead proviene de la comunicación y la sincronización:
\[
\boxed{\; T_o(p) = p\, T_{\text{comm}}. \;}
\]
Sustituyendo en la eficiencia se obtiene la forma canónica del curso:
\[
E = \frac{S_p}{p} = \frac{T_1}{p\,T_p} = \frac{W}{W + T_o}
  = \frac{1}{\,1 + \dfrac{T_o}{W}\,}.
\]
Para que $E=O(1)$ sea constante, el término de overhead $T_o/W$ debe permanecer constante: esa
igualdad, despejada, define la condición de isoeficiencia.

\subsubsection{Isoeficiencia en MPI (memoria distribuida)}

En MPI la comunicación es explícita y es el cuello de botella. Sustituyendo $n_{tr}=n_{te}=n$ en
$T_{\text{comm}}$ y conservando los términos dominantes:
\[
T_o = p\,T_{\text{comm}}
    \approx \underbrace{p^2\alpha}_{(1)\ \text{latencia}}
    \;+\; \underbrace{p\,\beta\, n\, d}_{(2)\ \text{scatter}}
    \;+\; \underbrace{p\,\beta\, n\, d\,\log p}_{(3)\ \text{bcast/reducción}}.
\]
El término (3) proviene de difundir el conjunto de prueba ($O(nd)$ bytes) y de reducir las listas
Top-$K$ a lo largo de una jerarquía de profundidad $\log p$; domina sobre (2). Imponemos
$T_o/W=\text{const}$ con $W=n^2 d$:
\[
(1)\quad \frac{p^2\alpha}{n^2 d}=\text{const} \;\Rightarrow\; n \propto p,
\qquad
(2)\quad \frac{p\,\beta n d}{n^2 d}=\frac{p\beta}{n}=\text{const} \;\Rightarrow\; n \propto p,
\]
\[
(3)\quad \frac{p\,\beta n d\,\log p}{n^2 d}=\frac{p\beta\log p}{n}=\text{const}
\;\Rightarrow\; n \propto p\,\log p.
\]
La condición más restrictiva es la del término (3). Por tanto:
\[
\boxed{\;
\text{Isoeficiencia MPI:}\quad n \propto p\,\log p,
\qquad
W_{iso} = \Theta\!\big(p^2 \log^2 p\big).
\;}
\]
Es decir, el dataset debe crecer ligeramente más rápido que $p$ (factor $\log p$) para amortizar
la difusión del test y la reducción global de candidatos sobre la red. Si la reducción se hiciera
con un \texttt{gather} plano al proceso raíz en vez de la reducción en árbol, el término dominante
sería $O(p)$ en lugar de $O(\log p)$, saturando el enlace del root y empeorando notablemente la
isoeficiencia.

\subsubsection{Isoeficiencia en OMP / PRAM (memoria compartida)}

En OpenMP los hilos comparten el mismo espacio de direcciones (modelo UMA), lo que cambia el
overhead de forma fundamental:
\begin{itemize}
    \item \textbf{No hay difusión de datos:} el conjunto de prueba es visible por todos los hilos,
    por lo que \emph{desaparecen} los términos $p\alpha$ y $\beta n d$ de la comunicación
    explícita ($T_{\text{comm}} \approx 0$).
    \item El único overhead proviene de la \textbf{sincronización de hilos} (barreras) y de la
    reducción de listas Top-$K$ de profundidad $O(\log p)$, realizada mediante accesos directos a
    memoria.
\end{itemize}
Modelando ese overhead (con $K$ constante):
\[
T_o^{\text{omp}} \approx p\cdot O(n_{te}K\log p) + O(p)
= O(p\,n\,\log p) + O(p).
\]
Imponiendo $T_o^{\text{omp}}/W=\text{const}$ con $W=n^2 d$, y observando que ya no aparece el
factor de transporte de datos $\beta d$:
\[
\frac{p\,n\,\log p}{n^2 d}=\frac{p\log p}{n\,d}=\text{const}
\;\Rightarrow\;
n \propto \frac{p\,\log p}{d}.
\]
\[
\boxed{\;
\text{Isoeficiencia OMP:}\quad n \propto \frac{p\,\log p}{d}\ \ (\text{a lo sumo } n\propto p),
\qquad
W_{iso} = \Theta(p^2)\ \text{sin la constante de red } \beta.
\;}
\]
OMP tiene una isoeficiencia asintóticamente mejor (no paga latencia ni ancho de banda de red),
\textbf{pero} $p$ está estrictamente acotado al número de núcleos físicos del nodo y todo el
dataset debe caber en su RAM.

\subsubsection{Comparación y elección del paradigma}

\begin{table}[H]
\centering
\begin{tabular}{lll}
\toprule
Criterio & MPI (memoria distribuida) & OMP (memoria compartida) \\
\midrule
Comunicación $T_{\text{comm}}$ & Explícita: $\alpha,\beta$ (red) & $\approx 0$ (acceso directo) \\
Overhead dominante & Difusión/reducción $\beta n d\log p$ & Sincronización (barreras) \\
Condición de isoeficiencia & $n \propto p\log p$ & $n \propto p\log p / d$ (más suave) \\
Trabajo $W_{iso}$ & $\Theta(p^2\log^2 p)$ & $\Theta(p^2)$ sin $\beta$ \\
Límite de $p$ & Horizontal (varios nodos) & Núcleos de un solo nodo \\
Límite de $n$ & Suma de RAM de los nodos & RAM de un solo nodo \\
\bottomrule
\end{tabular}
\caption{Comparación de la escalabilidad de ambos paradigmas para el KNN.}
\label{tab:iso-comparacion}
\end{table}

Aunque OMP escala mejor asintóticamente, su paralelismo está limitado a los núcleos y la memoria
de un solo nodo. En este proyecto se eligió \textbf{MPI por un tema de memoria}: al distribuir el
conjunto de entrenamiento, el mismo código puede superar el límite de RAM de un nodo y escalar
$n$ y $p$ horizontalmente sin reescritura. Los experimentos evaluaron hasta $p=32$ procesos, y la
elección del paradigma de memoria distribuida es la que habilita crecer más allá de ese punto
cuando el dataset no cabe en un único nodo.

